%!TeX encoding = UTF-8 Unicode
%!TeX spellcheck = en-UK
% !TeX root = Perrinet26mnesis.tex
% -*- coding: UTF-8; -*-
%%%-----------------------------------------------------------------
\begin{figure}[!ht]
  \begin{center}
    \includegraphics[width=\textwidth]{target.pdf}
  \end{center}
   \caption{Pattern retrieval by the recurrent HD-SNN after training. The network is clamped with input spikes (green crosses) comprising $50$ steps of spontaneous activity (left of the red dashed line) followed by the trigger window consisting of the initial sequence of one of the $M=16$ target motifs ($0 \leq t < D = 41$, delimited by the red and blue dashed lines). The motif shown, the second one, was chosen as characteristic of the average behaviour: with an intact trigger window, its $F_1$ score ($\approx 0.96$) is close to the median over the $M=16$ motifs ($\approx 0.95$). The autonomous replay is then simulated over the full sequence duration ($T = 1000$ steps). The raster plot shows network activity (blue plus signs, $N=1024$ neurons) which accurately matches the hidden target pattern which was cued during the trigger window. Note that the target spikes after the trigger window (red circles right of the blue dashed line) are hidden from the network.}\label{fig:target}
\end{figure}
%%%-----------------------------------------------------------------
